\documentclass{article}
\usepackage{apstats-poster}
% Formula IDs: ci-formula, margin-error, width-ci, random-condition, ten-pct-condition, normal-condition

\colorlet{PosterAccent}{UnitOrange}

\begin{document}
\begin{posterpage}

\PosterHeader[82pt]
  {UnitOrange}
  {3}
  {INFERENCE FOR PROPORTIONS \& CATEGORICAL DATA}
  {P09}
  {CONFIDENCE INTERVALS}
  {master recipe \quad•\quad also used in Units 4 and 5}
\InferenceSubstripes

\PosterZone{4.72in}{15.92in}{ZONE A · THE FORMULAS}{%
  \FormulaCardSized{44pt}{%
    \text{statistic}\ \pm\
    \underbrace{(\text{critical value})(\text{standard error})}_{\text{margin of error}}
  }
  \vspace{.14in}
  \begin{minipage}[t]{9.75in}
    \FormulaCardSized{43pt}{\mathrm{ME}=z^*\!\cdot\mathrm{SE}\quad(\text{proportions})}
  \end{minipage}%
  \hfill
  \begin{minipage}[t]{9.75in}
    \FormulaCardSized{43pt}{\mathrm{ME}=t^*\!\cdot\mathrm{SE}\quad(\text{means})}
    {\fontsize{21pt}{25pt}\selectfont Use the procedure's \(df\) to choose \(t^*\).}
  \end{minipage}%
  \vspace{.14in}
  \FormulaCardSized{41pt}{%
    \text{width}\propto\frac{1}{\sqrt{n}}
    \quad\Longrightarrow\quad
    \text{halve the margin }\Rightarrow\text{ quadruple }n
  }
  \vspace{.28in}
  \begin{minipage}[t]{12.65in}
    \MiniHeading{THE FOUR-STEP FRAME}\par\vspace{.12in}
    \StepCard{1}{STATE}{Name the parameter in context.}
    \StepCard{2}{PLAN}{Name the interval; check conditions.}
    \StepCard{3}{DO}{Formula, substitutions, and interval.}
    \StepCard{4}{CONCLUDE}{Interpret the interval in context.}
  \end{minipage}%
  \hfill
  \begin{minipage}[t]{6.75in}
    \MiniHeading{COMMON \(z^*\) VALUES}\par\vspace{.12in}
    \TextCard{%
      \renewcommand{\arraystretch}{1.35}
      \centering
      \begin{tabular}{>{\bfseries}c c}
        confidence & \(z^*\) \\
        \midrule
        90\% & 1.645 \\
        95\% & 1.960 \\
        99\% & 2.576
      \end{tabular}
    }
    \vspace{.18in}
    {\fontsize{24pt}{30pt}\selectfont
      Get \(t^*\) from the table or calculator using \textbf{df}.}
  \end{minipage}%
}

\PosterZone{16.07in}{20.18in}{ZONE B · CONDITIONS — ALWAYS IN THIS ORDER}{%
  \begin{minipage}[t]{6.15in}
    \MiniHeading{1 · RANDOM}\par\vspace{.08in}
    {\fontsize{25pt}{31pt}\selectfont
      Random sample or\par randomized experiment.}
  \end{minipage}%
  \hfill
  \begin{minipage}[t]{6.15in}
    \MiniHeading{2 · 10\%}\par\vspace{.08in}
    {\fontsize{25pt}{31pt}\selectfont
      Sampling without replacement:\par
      \(n<0.10N\).}
  \end{minipage}%
  \hfill
  \begin{minipage}[t]{6.9in}
    \MiniHeading{3 · NORMAL / LARGE COUNTS}\par\vspace{.08in}
    {\fontsize{23pt}{29pt}\selectfont
      CI: \(n\hat p\ge10\), \(n(1-\hat p)\ge10\)\par
      Test: \(np_0\ge10\), \(n(1-p_0)\ge10\) (see P11)\par
      means: normal population, \(n\ge30\), or no strong skew/outliers}
  \end{minipage}%
}

\PosterZone{20.33in}{23.62in}{ZONE C · SAY IT IN WORDS}{%
  \SentenceFrame{“We are \PosterBlank{1.0in}\% confident that the interval
    from \PosterBlank{1.25in} to \PosterBlank{1.25in} captures the true
    \PosterBlank{4.1in}.”}
  \vspace{.12in}
  {\fontsize{24pt}{30pt}\selectfont
    \textbf{Confidence level:} If we took many samples and built an interval
    each time,\par
    about \PosterBlank{.8in}\% of those intervals would capture the true
    \PosterBlank{3.0in}.}
}

\PosterZone{23.77in}{27.45in}{ZONE D · WATCH OUT}{%
  \begin{minipage}[t]{19.15in}
    {\posterheading\bfseries\color{PosterAccent}\fontsize{31pt}{36pt}\selectfont
      THE INTERVAL ESTIMATES THE PARAMETER. THE LEVEL DESCRIBES THE METHOD.}\par
    \vspace{.12in}
    {\fontsize{29pt}{37pt}\selectfont
      Never say the fixed parameter has a 95\% chance of being inside the
      interval. After the interval is calculated, it either captures the
      parameter or it does not.}
  \end{minipage}%
}

\WorksheetTag{u6\_l2–3 · u6\_l8–9 · u7\_l2–3 · u7\_l6–7}

\end{posterpage}
\end{document}
